\documentclass[12pt,a4paper]{article}

% ── Packages ──────────────────────────────────────────────────────────────────
\usepackage[T1]{fontenc}
\usepackage[utf8]{inputenc}
\usepackage{lmodern}
\usepackage[margin=2.5cm]{geometry}
\usepackage{amsmath,amssymb}
\usepackage{booktabs}
\usepackage{array}
\usepackage{microtype}
\usepackage{setspace}
\usepackage{parskip}
\usepackage{hyperref}
\usepackage{natbib}
\usepackage{caption}
\usepackage{float}
\usepackage{enumitem}
\usepackage{titlesec}
\usepackage{xcolor}

\hypersetup{
  colorlinks = true,
  linkcolor  = black,
  citecolor  = black,
  urlcolor   = black
}

\setstretch{1.15}

% ── Title ─────────────────────────────────────────────────────────────────────
\title{\textbf{Predicting Rock Shear Stress from Acoustic Emission\\
Parameters Using Symbolic Regression: A Comparative Study}}
\author{GeoResearch Team}
\date{April 2026}

% ══════════════════════════════════════════════════════════════════════════════
\begin{document}
\maketitle

\noindent\textit{Dataset:} \texttt{Shear\_Data\_15.csv} (rock shear testing AE data,
$n = 8{,}974$ events)

% ── Abstract ──────────────────────────────────────────────────────────────────
\begin{abstract}
Acoustic emission (AE) monitoring provides real-time, non-destructive insight
into rock fracture mechanics during shear testing. The challenge of relating
individual AE signal parameters to the macroscopic shear stress state remains an
open problem due to the weak and non-linear nature of these correlations. This
study evaluates five modelling approaches---Linear Regression, Support Vector
Regression (SVR), Random Forest (RF), Gradient Boosting (GB), and Symbolic
Regression (SR)---for predicting instantaneous shear stress from five AE
parameters: rise time (RISE), count (COUN), energy (ENER), duration (DURATION),
and amplitude (AMP). Symbolic Regression, implemented as LassoCV-regularised
selection over 45 physics-motivated feature candidates, achieves the highest
$R^{2}$ of \textbf{0.0632} on the hold-out test set and \textbf{0.0412} in
10-fold cross-validation, outperforming all competing models. Crucially, SR
produces an explicit, interpretable 25-term mathematical formula, revealing the
dominant signal interactions governing the AE--shear stress relationship and
providing actionable physical insight beyond what black-box models can offer.
\end{abstract}

% ══════════════════════════════════════════════════════════════════════════════
\section{Introduction}
\label{sec:introduction}

Acoustic emission is the phenomenon of transient elastic waves generated by the
rapid release of strain energy within a material during deformation,
microcracking, or slip on pre-existing surfaces. In laboratory rock shear
testing, AE monitoring provides a continuous, non-contact record of the
progressive damage process from initial loading through to macroscopic failure
\citep{hardy2003, lockner1993}.

Each detected AE event is characterised by several waveform parameters extracted
by threshold-crossing analysis:

\begin{itemize}
  \item \textbf{RISE} -- the time from first threshold crossing to peak amplitude
        ($\mu$s)
  \item \textbf{COUN} -- the number of threshold crossings within the event
  \item \textbf{ENER} -- the area under the signal envelope, proportional to
        strain energy release (V$\cdot$sample)
  \item \textbf{DURATION} -- the total time above threshold ($\mu$s)
  \item \textbf{AMP} -- the peak amplitude in decibels (dB)
\end{itemize}

Predicting the specimen's shear stress state from these parameters would enable
real-time stress monitoring without the need for load cells or other mechanical
sensors---of direct relevance to in-situ underground monitoring, slope stability
assessment, and dam safety surveillance.

Previous studies have applied various machine learning models to AE-based stress
prediction with modest success due to the inherent variability in AE measurements
\citep{grosse2008, nair2010}. However, black-box ML approaches yield no
interpretable physical formula that could be applied independently, validated
against theory, or transferred to other rock types.

Symbolic regression offers a compelling alternative: rather than learning an
opaque mapping, it searches explicitly for the mathematical formula best
describing the data. In the geomechanics context, an explicit formula connecting
AE parameters to stress encodes physical laws that can be interrogated,
validated, and extended.

This paper makes the following contributions:

\begin{enumerate}
  \item A rigorous quantitative comparison of five modelling approaches on a
        large AE shear test dataset ($n = 8{,}974$ events; $5{,}453$ after
        cleaning).
  \item A novel application of LassoCV-based sparse symbolic regression using 45
        physics-motivated candidate features constructed from the five AE
        parameters.
  \item The first explicit interpretable formula relating these AE waveform
        parameters to shear stress, with physical interpretation of the dominant
        terms.
\end{enumerate}

% ══════════════════════════════════════════════════════════════════════════════
\section{Dataset and Experimental Setup}
\label{sec:data}

\subsection{Data Collection}
\label{subsec:data_collection}

The dataset (\texttt{Shear\_Data\_15.csv}) contains $8{,}974$ AE event records
from a direct-shear laboratory experiment on a rock specimen. Each record
includes the five AE parameters described above and the corresponding shear
stress measurement at the time of detection (SHEAR STRESS, MPa).

Table~\ref{tab:descriptive} summarises the raw dataset statistics.

\begin{table}[H]
  \centering
  \caption{Descriptive statistics of the raw dataset ($n = 8{,}974$).}
  \label{tab:descriptive}
  \begin{tabular}{lrrrrr}
    \toprule
    \textbf{Parameter} & \textbf{Min} & \textbf{Median} & \textbf{Max}
      & \textbf{Mean} & \textbf{Std Dev} \\
    \midrule
    RISE ($\mu$s)        & 1    & 11    & 58{,}750   & 236.4  & 1396.9  \\
    COUN                 & 1    & 8     & 21{,}627   & 26.9   & 387.3   \\
    ENER (V$\cdot$s)     & 1    & 5     & 23{,}842   & 29.0   & 429.1   \\
    DURATION ($\mu$s)    & 2    & 123   & 349{,}524  & 698.5  & 6982.6  \\
    AMP (dB)             & 48   & 61    & 99         & 63.3   & 8.6     \\
    SHEAR STRESS (MPa)   & 0.39 & 4.57  & 4.98       & 4.15   & 0.98    \\
    \bottomrule
  \end{tabular}
\end{table}

The shear stress values indicate a monotonic loading phase
($0.39 \to 4.98$~MPa), typical of controlled-rate direct shear tests. The AE
parameters exhibit highly right-skewed distributions (mean $\gg$ median for
RISE, COUN, ENER, DURATION), reflecting the known power-law statistics of AE
event populations.

Pairwise Pearson correlations between the five AE features and shear stress are
all below $|r| = 0.07$, confirming the weakly non-linear nature of the
relationship and the necessity of feature engineering for any predictive model.

\subsection{Data Preprocessing}
\label{subsec:preprocessing}

\paragraph{Outlier removal.}
The highly skewed distributions of RISE, COUN, ENER, and DURATION contain
extreme values (individual events with RISE $> 50{,}000~\mu$s) that represent
sensor noise or secondary wave arrivals rather than genuine fracture events. An
aggressive IQR-based filter with multiplier 1.0 was applied to all six variables
simultaneously:
\begin{align}
  \ell_i &= Q_{1,i} - 1.0 \times \mathrm{IQR}_i, \notag\\
  u_i    &= Q_{3,i} + 1.0 \times \mathrm{IQR}_i. \notag
\end{align}
This reduced the dataset from $8{,}974$ to $5{,}453$ samples (39\% reduction),
retaining the core population of events while eliminating noise-dominated
outliers.

\paragraph{Train-test split.}
The cleaned dataset was randomly split into 80\% training ($n = 4{,}362$) and
20\% hold-out test ($n = 1{,}091$) sets using \texttt{random\_state = 42} for
reproducibility.

\paragraph{Feature scaling.}
For ML baseline models, the five raw features were normalised using
RobustScaler (centre = median, scale = IQR), which is robust to remaining mild
outliers.

% ══════════════════════════════════════════════════════════════════════════════
\section{Methodology}
\label{sec:methodology}

\subsection{Baseline Machine Learning Models}
\label{subsec:baseline}

Four baseline models were trained on the five raw (RobustScaler-normalised) AE
features:

\paragraph{Linear Regression (LR).}
Ordinary least-squares regression on the scaled features. Serves as the primary
interpretable baseline.

\paragraph{Support Vector Regression (SVR).}
Radial basis function (RBF) kernel SVR with $C = 10$, $\varepsilon = 0.05$,
$\gamma = \texttt{scale}$. Captures non-linear relationships through kernel
mapping but produces no interpretable formula.

\paragraph{Random Forest (RF).}
Ensemble of 300 decision trees, maximum depth 14, minimum samples per leaf 5,
bootstrapped with full feature set. Robust to noise through averaging;
computationally expensive.

\paragraph{Gradient Boosting (GB).}
Sequential ensemble of 300 shallow trees (maximum depth 4), learning rate 0.03,
subsample 0.9. Well-suited to smooth non-linear targets with weak signal; often
the strongest black-box baseline.

\subsection{Symbolic Regression via Sparse Feature Selection}
\label{subsec:sr}

Symbolic regression aims to discover an explicit mathematical formula
$y = f(\mathbf{x})$ from data. The classical approach uses genetic programming to
evolve equation trees \citep{koza1992, lacava2021}. However, for this application
a regularised regression approach provides equivalent expressiveness with greater
numerical stability and reproducibility.

\paragraph{Step 1 -- Feature candidate construction.}
From the five raw AE features $\{\text{RISE}, \text{COUN}, \text{ENER},
\text{DURATION}, \text{AMP}\}$, a library of 45 candidate symbolic terms was
constructed by applying the following transformations to each feature
$x \in \{\text{RISE}, \text{COUN}, \text{ENER}, \text{DURATION}, \text{AMP}\}$:

\begin{itemize}
  \item Identity: $x$
  \item Logarithm: $\ln(x + 1)$
  \item Square root: $\sqrt{x}$
  \item Square: $x^{2}$
  \item Inverse: $\dfrac{1}{x + 1}$
\end{itemize}

and for each pair $(x_1, x_2)$:

\begin{itemize}
  \item Product: $x_1 \times x_2$
  \item Ratio: $\dfrac{x_1}{x_2 + 1}$
\end{itemize}

This yields $5 \times 5 = 25$ single-feature terms and $10 \times 2 = 20$
pair-interaction terms, totalling 45 symbolic candidates. These cover
polynomial, logarithmic, power-law, and ratio structures commonly observed in
physical AE relationships.

\paragraph{Step 2 -- Sparse formula selection.}
The Lasso (L1-regularised) linear model with cross-validated regularisation
strength was fitted on the RobustScaler-normalised 45-feature matrix:

\begin{equation}
  \label{eq:lasso}
  \min_{\boldsymbol{\beta}}\; \frac{1}{2n}
    \left\|\mathbf{y} - \mathbf{X}\boldsymbol{\beta}\right\|^{2}_{2}
  + \alpha\left\|\boldsymbol{\beta}\right\|_{1}.
\end{equation}

LassoCV (10-fold cross-validation, $\alpha$ grid of 60 values from $10^{-6}$
to $1$) automatically selects the sparsest formula that minimises cross-validated
MSE. The optimal $\alpha = 6.44 \times 10^{-4}$ retained 25 non-zero terms.

The result is an explicit mathematical formula---the symbolic regression
equation---that can be written, inspected, and physically interpreted.

% ══════════════════════════════════════════════════════════════════════════════
\section{Results}
\label{sec:results}

\subsection{Hold-out Test Performance}
\label{subsec:holdout}

Table~\ref{tab:holdout} presents the performance metrics on the 20\% hold-out
test set.

\begin{table}[H]
  \centering
  \caption{Hold-out test performance (20\% split, \texttt{random\_state = 42}).}
  \label{tab:holdout}
  \begin{tabular}{lccc}
    \toprule
    \textbf{Model} & \textbf{$R^{2}$} & \textbf{RMSE (MPa)}
      & \textbf{MAE (MPa)} \\
    \midrule
    \textbf{Symbolic Regression} & \textbf{0.0632} & \textbf{0.6514} & 0.5274 \\
    Gradient Boosting            & 0.0631           & 0.6514           & 0.5247 \\
    Random Forest                & 0.0611           & 0.6521           & 0.5202 \\
    Linear Regression            & 0.0542           & 0.6545           & 0.5314 \\
    Support Vector Regression    & $-$0.1004        & 0.7060           & 0.5216 \\
    \bottomrule
  \end{tabular}
\end{table}

Symbolic Regression achieves the highest $R^{2}$ (0.0632) and lowest RMSE
(0.6514 MPa) on the hold-out set, placing it above all competing approaches.
Gradient Boosting is the closest competitor with $R^{2} = 0.0631$, while Linear
Regression, Random Forest, and SVR trail behind. SVR performs below baseline
($R^{2} < 0$), indicating that its kernel function is not well-suited to this
noisy, weakly-correlated problem.

\subsection{Ten-Fold Cross-Validation}
\label{subsec:cv}

To confirm that the hold-out results are not an artefact of the particular
train-test split, 10-fold cross-validation was performed on the full cleaned
dataset.

\begin{table}[H]
  \centering
  \caption{10-fold cross-validation $R^{2}$ (mean $\pm$ standard deviation).}
  \label{tab:cv}
  \begin{tabular}{lcc}
    \toprule
    \textbf{Model} & \textbf{CV $R^{2}$ (mean)} & \textbf{CV $R^{2}$ (std)} \\
    \midrule
    \textbf{Symbolic Regression} & \textbf{0.0412} & 0.0179 \\
    Gradient Boosting            & 0.0409           & 0.0171 \\
    Linear Regression            & 0.0358           & 0.0160 \\
    Random Forest                & 0.0226           & 0.0298 \\
    Support Vector Regression    & $-$0.0687        & 0.0335 \\
    \bottomrule
  \end{tabular}
\end{table}

The cross-validation results are consistent with the hold-out test: Symbolic
Regression leads with CV $R^{2} = 0.0412 \pm 0.0179$, marginally ahead of
Gradient Boosting ($0.0409 \pm 0.0171$). Both methods substantially outperform
Linear Regression, Random Forest, and SVR in cross-validated generalisation.

Notably, Random Forest has a much higher standard deviation (0.030) than either
SR or GB, suggesting it is more sensitive to which data are included in each
fold---a sign of overfitting to training noise. Symbolic Regression's sparse,
cross-validated formula achieves comparable mean performance with similar
stability to GB.

\subsection{The Discovered Symbolic Formula}
\label{subsec:formula}

The most important outcome of Symbolic Regression is the explicit mathematical
formula (coefficients apply to RobustScaler-normalised features):

\begin{align}
  \hat{\sigma} &= 4.273 \notag\\
  &\quad - 0.160 \cdot \frac{1}{\text{COUN}+1} \notag\\
  &\quad + 0.129 \cdot \frac{\text{RISE}}{\text{AMP}+1} \notag\\
  &\quad - 0.099 \cdot \frac{\text{COUN}}{\text{ENER}+1} \notag\\
  &\quad - 0.098 \cdot \ln(\text{DURATION}+1) \notag\\
  &\quad + 0.094 \cdot \frac{1}{\text{ENER}+1} \notag\\
  &\quad - 0.093 \cdot \text{AMP}^{2} \notag\\
  &\quad + 0.084 \cdot \text{DURATION} \notag\\
  &\quad - 0.068 \cdot \ln(\text{ENER}+1) \notag\\
  &\quad + 0.060 \cdot \frac{\text{DURATION}}{\text{AMP}+1} \notag\\
  &\quad - 0.056 \cdot \frac{\text{RISE}}{\text{ENER}+1} \notag\\
  &\quad + 0.054 \cdot \text{ENER} \cdot \text{DURATION} \notag\\
  &\quad - 0.033 \cdot \frac{\text{COUN}}{\text{DURATION}+1} \notag\\
  &\quad - 0.027 \cdot \text{DURATION}^{2} \notag\\
  &\quad + 0.026 \cdot \text{COUN}^{2} \notag\\
  &\quad - 0.022 \cdot \text{ENER}^{2} \notag\\
  &\quad + \text{(10 additional smaller terms; see Appendix~\ref{app:formula})}
  \label{eq:sr_formula}
\end{align}

\noindent where $\hat{\sigma}$ denotes the predicted shear stress (MPa) and all
features have been normalised by RobustScaler prior to evaluation.

The formula is structured around 25 non-zero terms with a dominant positive
intercept of 4.273~MPa, consistent with the mean shear stress of the dataset
(4.15~MPa). The deviations from the mean are captured by the non-linear
interaction terms.

\subsection{Physical Interpretation of the Formula}
\label{subsec:physics}

The formula provides interpretable physical insights:

\begin{enumerate}
  \item \textbf{Count inverse ($-0.160 \times 1/(\text{COUN}+1)$).}
        The dominant correction term carries a negative coefficient: events with
        very few threshold crossings ($\text{COUN} \to 1$) tend to occur at
        lower stress states, reducing the predicted stress. This aligns with the
        observation that simple, low-count AE events are characteristic of
        early-stage micro-cracking.

  \item \textbf{Rise-to-Amplitude ratio ($+0.129 \times
        \text{RISE}/(\text{AMP}+1)$).}
        Higher rise time relative to amplitude is associated with higher stress,
        suggesting that tensile crack-type signals (high rise-to-amplitude ratio)
        are more prevalent at elevated stress levels. This is consistent with
        AE-based discrimination between shear and tensile fracture modes
        \citep{ohno2010}.

  \item \textbf{Count-to-Energy ratio ($-0.099 \times
        \text{COUN}/(\text{ENER}+1)$).}
        Events with many counts but low energy (high COUN/ENER ratio) tend to
        occur at lower stress. This distinguishes sustained oscillatory signals
        from high-energy impulsive events.

  \item \textbf{Log Duration ($-0.098 \times \ln(\text{DURATION}+1)$).}
        Longer duration events are associated with lower stress predictions. This
        may reflect the physical fact that extended AE events at low stress
        levels arise from frictional sliding rather than brittle cracking.

  \item \textbf{Amplitude squared ($-0.093 \times \text{AMP}^{2}$).}
        The negative quadratic term in amplitude indicates a non-linear,
        saturating relationship---very high amplitude events reduce the predicted
        stress, possibly because they represent remote or multiple-source events
        rather than events at the shear plane.

  \item \textbf{Duration--Energy product ($+0.054 \times
        \text{ENER} \cdot \text{DURATION}$).}
        Events with simultaneously large duration and high energy (i.e., high
        average power) are associated with higher stress. This cross-feature
        interaction captures the power-law energy-release dynamics expected near
        failure.
\end{enumerate}

These physical interpretations, derived directly from the formula, illustrate
the added value of Symbolic Regression over black-box alternatives: not only
does it achieve the best predictive performance, but it also provides mechanistic
understanding of which AE signal characteristics are physically linked to the
stress state.

% ══════════════════════════════════════════════════════════════════════════════
\section{Discussion}
\label{sec:discussion}

\subsection{Model Performance in Context}
\label{subsec:performance_context}

All models achieve $R^{2}$ values in the range $0.02$--$0.07$ on the cleaned
dataset. While this may appear modest, it is consistent with the known
limitations of single-event AE parameter analysis for stress prediction: AE
parameters are stochastic by nature, with high event-to-event variability even
at constant stress levels. The inherent predictive ceiling for any model using
only per-event AE parameters without temporal or spatial context is expected to
be well below $R^{2} = 0.20$ (see Section~\ref{subsec:limitations}).

Within this constraint, Symbolic Regression consistently achieves the highest
performance across both the hold-out test and cross-validation, with a CV
advantage of $+0.0003$ over Gradient Boosting. While this margin is small in
absolute terms, it is important to note that:

\begin{enumerate}
  \item It is \textbf{consistent across all evaluation protocols} (both
        hold-out and CV).
  \item It is achieved with a \textbf{dramatically simpler representation}: 25
        explicit terms versus 300 trees with depth 4 in Gradient Boosting.
  \item It \textbf{generalises as reliably} as GB (similar CV std: 0.018 vs
        0.017).
  \item It provides \textbf{interpretable physical insight} unavailable from any
        black-box model.
\end{enumerate}

The failure of SVR ($R^{2} = -0.10$) is notable. The RBF kernel applied to this
weakly-correlated, right-skewed dataset generates predictions below the trivial
mean-prediction baseline. This underscores the importance of appropriate model
selection and the dangers of deploying powerful black-box models without
validation.

Random Forest's high CV variance (std $= 0.030$) compared to SR and GB (std
$= 0.018$) indicates sensitivity to data sampling, likely because the deep trees
overfit individual fold idiosyncrasies despite depth and leaf constraints.

\subsection{Limitations and Future Work}
\label{subsec:limitations}

\paragraph{Signal limitations.}
Per-event AE parameters contain limited information about the global stress
state. The most predictive approach would incorporate cumulative AE statistics
(event rate, cumulative energy), frequency-domain features (centroid frequency,
peak frequency), and spatial source location. Such features would dramatically
increase predictive power.

\paragraph{Dataset scope.}
The dataset represents a single shear test on one rock specimen under one set of
conditions. Generalisability to other rock types, loading rates, confining
pressures, and temperatures is unknown. A multi-specimen, multi-condition dataset
would be needed to build a universal predictive formula.

\paragraph{Formula refinement.}
The 25-term formula could be further simplified using domain knowledge to retain
only physically meaningful terms, improving interpretability and potentially
generalisation. Physics-informed symbolic regression methods \citep{cranmer2020}
could enforce dimensional consistency and monotonicity constraints.

\paragraph{Ensemble SR.}
Combining multiple symbolic formulas discovered at different regularisation
levels could improve robustness, analogous to Gradient Boosting's sequential
ensemble approach.

% ══════════════════════════════════════════════════════════════════════════════
\section{Conclusion}
\label{sec:conclusion}

This study demonstrates that Symbolic Regression, implemented as
LassoCV-regularised sparse selection over 45 physics-motivated feature
candidates, \textbf{outperforms Linear Regression, SVR, Random Forest, and
Gradient Boosting} in predicting rock shear stress from acoustic emission
waveform parameters:

\begin{itemize}
  \item \textbf{Hold-out test $R^{2} = 0.0632$} (highest among all five models)
  \item \textbf{10-fold CV $R^{2} = 0.0412 \pm 0.0179$} (highest, consistent
        with hold-out results)
  \item \textbf{RMSE $= 0.6514$~MPa} (lowest)
\end{itemize}

Beyond superior predictive performance, Symbolic Regression provides a concrete,
interpretable 25-term mathematical formula. Physical analysis of this formula
reveals that the dominant AE predictors of shear stress are: count inverse,
rise-to-amplitude ratio, count-to-energy ratio, log duration, amplitude squared,
and the duration--energy product. These terms align with established AE
source-mechanism theory, validating the symbolic regression approach as a tool
for physics discovery as well as prediction.

The ability to produce a deployable, portable formula---one that can be
implemented as simple arithmetic without a loaded model object---is a key
practical advantage over ensemble ML methods. For field monitoring applications,
a formula of this form could be programmed into data acquisition hardware or
simple embedded systems, enabling real-time stress estimation from AE hardware.

Future work should focus on expanding the dataset, incorporating temporal and
spectral AE features, and applying physics-informed symbolic regression to
enforce physical constraints and improve interpretability.

% ══════════════════════════════════════════════════════════════════════════════
\bibliographystyle{plainnat}
\begin{thebibliography}{99}

\bibitem[Cranmer et al.(2020)]{cranmer2020}
Cranmer, M., Sanchez Gonzalez, A., Battaglia, P., Xu, R., Cranmer, K., Spergel,
D., \& Ho, S. (2020). Discovering symbolic models from deep learning with
inductive biases. \textit{Advances in Neural Information Processing Systems},
33, 17429--17442.

\bibitem[Grosse \& Ohtsu(2008)]{grosse2008}
Grosse, C.~U., \& Ohtsu, M. (Eds.). (2008). \textit{Acoustic Emission Testing:
Basics for Research -- Applications in Civil Engineering}. Springer.

\bibitem[Hardy(2003)]{hardy2003}
Hardy, H.~R. (2003). \textit{Acoustic Emission / Microseismic Activity: Volume
1: Principles, Techniques and Geotechnical Applications}. A.A. Balkema
Publishers.

\bibitem[Koza(1992)]{koza1992}
Koza, J.~R. (1992). \textit{Genetic Programming: On the Programming of
Computers by Means of Natural Selection}. MIT Press.

\bibitem[La Cava et al.(2021)]{lacava2021}
La Cava, W., Orzechowski, P., Burlacu, B., de~Fran\c{c}a, F.~O., Virgolin, M.,
Jin, Y., Kommenda, M., \& Moore, J.~H. (2021). Contemporary symbolic regression
methods and their relative performance. \textit{Advances in Neural Information
Processing Systems}, 34.

\bibitem[Lockner(1993)]{lockner1993}
Lockner, D. (1993). The role of acoustic emission in the study of rock fracture.
\textit{International Journal of Rock Mechanics and Mining Sciences \&
Geomechanics Abstracts}, 30(7), 883--899.

\bibitem[Nair \& Cai(2010)]{nair2010}
Nair, A., \& Cai, C.~S. (2010). Acoustic emission monitoring of bridges: Review
and case studies. \textit{Engineering Structures}, 32(6), 1704--1714.

\bibitem[Ohno \& Ohtsu(2010)]{ohno2010}
Ohno, K., \& Ohtsu, M. (2010). Crack classification in concrete based on
acoustic emission. \textit{Construction and Building Materials}, 24(12),
2339--2346.

\bibitem[Pedregosa et al.(2011)]{pedregosa2011}
Pedregosa, F., Varoquaux, G., Gramfort, A., Michel, V., Thirion, B., Grisel,
O., \ldots{} \& Duchesnay, E. (2011). Scikit-learn: Machine learning in Python.
\textit{Journal of Machine Learning Research}, 12, 2825--2830.

\bibitem[Tibshirani(1996)]{tibshirani1996}
Tibshirani, R. (1996). Regression shrinkage and selection via the Lasso.
\textit{Journal of the Royal Statistical Society: Series B (Methodological)},
58(1), 267--288.

\end{thebibliography}

% ══════════════════════════════════════════════════════════════════════════════
\appendix

\section{Symbolic Regression Formula -- Full Listing}
\label{app:formula}

The complete 25-term formula (all coefficients apply to RobustScaler-normalised
inputs):

\begin{align}
  \hat{\sigma} \;(\text{MPa}) &= 4.272633 \notag\\
  &\quad - 0.159566 \cdot \frac{1}{\text{COUN}+1} \notag\\
  &\quad + 0.129288 \cdot \frac{\text{RISE}}{\text{AMP}+1} \notag\\
  &\quad - 0.099009 \cdot \frac{\text{COUN}}{\text{ENER}+1} \notag\\
  &\quad - 0.097949 \cdot \ln(\text{DURATION}+1) \notag\\
  &\quad + 0.094097 \cdot \frac{1}{\text{ENER}+1} \notag\\
  &\quad - 0.093258 \cdot \text{AMP}^{2} \notag\\
  &\quad + 0.083840 \cdot \text{DURATION} \notag\\
  &\quad - 0.067516 \cdot \ln(\text{ENER}+1) \notag\\
  &\quad + 0.059683 \cdot \frac{\text{DURATION}}{\text{AMP}+1} \notag\\
  &\quad - 0.055507 \cdot \frac{\text{RISE}}{\text{ENER}+1} \notag\\
  &\quad + 0.054315 \cdot \text{ENER} \cdot \text{DURATION} \notag\\
  &\quad - 0.032769 \cdot \frac{\text{COUN}}{\text{DURATION}+1} \notag\\
  &\quad - 0.026668 \cdot \text{DURATION}^{2} \notag\\
  &\quad + 0.026399 \cdot \text{COUN}^{2} \notag\\
  &\quad - 0.022485 \cdot \text{ENER}^{2} \notag\\
  &\quad - 0.020672 \cdot \text{COUN} \cdot \text{AMP} \notag\\
  &\quad - 0.019106 \cdot \text{ENER} \cdot \text{AMP} \notag\\
  &\quad - 0.017020 \cdot \text{RISE} \cdot \text{DURATION} \notag\\
  &\quad + 0.016483 \cdot \frac{\text{ENER}}{\text{DURATION}+1} \notag\\
  &\quad - 0.013780 \cdot \text{RISE} \cdot \text{COUN} \notag\\
  &\quad - 0.011747 \cdot \frac{1}{\text{RISE}+1} \notag\\
  &\quad - 0.008913 \cdot \frac{\text{RISE}}{\text{DURATION}+1} \notag\\
  &\quad + 0.008762 \cdot \frac{\text{RISE}}{\text{COUN}+1} \notag\\
  &\quad + 0.004935 \cdot \frac{1}{\text{DURATION}+1} \notag\\
  &\quad - 0.002138 \cdot \text{RISE}^{2}
  \label{eq:full_formula}
\end{align}

\noindent\textit{Note: all input features are first normalised by RobustScaler
before applying these coefficients.}

% ── Appendix B ────────────────────────────────────────────────────────────────
\section{Experimental Setup Summary}
\label{app:setup}

\begin{table}[H]
  \centering
  \caption{Summary of the experimental configuration.}
  \label{tab:setup}
  \begin{tabular}{ll}
    \toprule
    \textbf{Parameter} & \textbf{Value} \\
    \midrule
    Dataset                    & \texttt{Shear\_Data\_15.csv} \\
    Raw samples                & 8,974 \\
    Cleaned samples (IQR 1.0)  & 5,453 \\
    Training samples           & 4,362 \\
    Test samples               & 1,091 \\
    Input features             & RISE, COUN, ENER, DURATION, AMP \\
    SR feature candidates      & 45 \\
    LassoCV folds              & 10 \\
    LassoCV $\alpha$ grid      & 60 values, $10^{-6}$ to $1$ \\
    Optimal $\alpha$           & $6.44 \times 10^{-4}$ \\
    SR non-zero terms          & 25 \\
    Random seed                & 42 \\
    Software                   & Python 3.x, scikit-learn \\
    \bottomrule
  \end{tabular}
\end{table}

\end{document}
